\section{Results}
\label{sec:results}
\subsection{Sealed efficacy, controls, and additional seeds}
% AUTO-GENERATED by analyze_study2_secondary.py. DO NOT EDIT.
\begin{table*}[t]
\centering
\caption{Headline sealed Study 2 comparison over the same
\studyclaim{StudyTwoDesignCount} designs. Penalized
Fmax is averaged equally over designs and scores incorrect or implementation-
failed samples as zero. Confidence intervals are paired, stratified 95\%
bootstrap intervals versus the SFT baseline.}
\label{tab:study2-primary}
\small
\begin{tabular}{lrrrr}
\toprule
Arm & Training seeds & Penalized Fmax [MHz] & Correct [\%] & $\Delta$ vs. SFT [95\% CI] \\
\midrule
SFT & -- & \studyclaim{StudyTwoOverallSftPenalized} &
\studyclaim{StudyTwoSftCorrectness} & -- \\
Correctness only & \studyclaim{StudyTwoCorrectnessTrainingSeeds} &
\studyclaim{StudyTwoOverallCorrectnessOnlyPenalized} &
\studyclaim{StudyTwoCorrectnessOnlyCorrectness} &
\studyclaim{StudyTwoOverallCorrectnessOnlySftGain}
[\studyclaim{StudyTwoOverallCorrectnessOnlySftCiLow},
 \studyclaim{StudyTwoOverallCorrectnessOnlySftCiHigh}] \\
Original MLP reward & \studyclaim{StudyTwoMlpTrainingSeeds} &
\studyclaim{StudyTwoOverallMlpPenalized} & \studyclaim{StudyTwoMlpCorrectness} &
\studyclaim{StudyTwoOverallMlpSftGain}
[\studyclaim{StudyTwoOverallMlpSftCiLow},
 \studyclaim{StudyTwoOverallMlpSftCiHigh}] \\
Repaired RF reward & \studyclaim{StudyTwoRfTrainingSeeds} &
\studyclaim{StudyTwoOverallRfPenalized} & \studyclaim{StudyTwoRfCorrectness} &
\studyclaim{StudyTwoOverallRfSftGain}
[\studyclaim{StudyTwoOverallRfSftCiLow},
 \studyclaim{StudyTwoOverallRfSftCiHigh}] \\
\bottomrule
\end{tabular}
\end{table*}


The repaired-RF policy improves the matched physical endpoint across the
\studyclaim{StudyTwoDesignCount} held-out designs
(Table~\ref{tab:study2-primary}). Penalized equal-sample \fmax{} rises from
\studyclaim{StudyTwoOverallSftPenalized} to
\studyclaim{StudyTwoOverallRfPenalized}~MHz.  The paired gain is
\studyclaim{StudyTwoOverallRfSftGain}~MHz with interval
[\studyclaim{StudyTwoOverallRfSftCiLow},
\studyclaim{StudyTwoOverallRfSftCiHigh}]~MHz. The gain appears in both regimes:
interpolation improves by
\studyclaim{StudyTwoInterpRfSftGain}~MHz with interval
[\studyclaim{StudyTwoInterpRfSftCiLow},
\studyclaim{StudyTwoInterpRfSftCiHigh}], while extrapolation improves by
\studyclaim{StudyTwoExtrapRfSftGain}~MHz with interval
[\studyclaim{StudyTwoExtrapRfSftCiLow},
\studyclaim{StudyTwoExtrapRfSftCiHigh}]~MHz.
Fig.~\ref{fig:study2-primary} shows every design together with the frozen paired
analysis. Zero-valued designs remain in the figure and in the aggregate, so
the reported improvement accounts for failures as well as successful hardware.

\begin{figure*}[t]
\centering
\includegraphics[width=.82\textwidth]{fig02_study2_primary}
\caption{Sealed primary efficacy.  Panel (a) retains every design and joins its
SFT and repaired-RF penalized equal-sample post-route \fmax.  Panel (b) reports
the frozen paired effects and the separately displayed training-seed effects.
The overall gain is \studyclaim{StudyTwoOverallRfSftGain}~MHz
[\studyclaim{StudyTwoOverallRfSftCiLow},
\studyclaim{StudyTwoOverallRfSftCiHigh}]~MHz.}
\label{fig:study2-primary}
\end{figure*}

Both independently trained RF policies improve, with overall
effects of \studyclaim{StudyTwoRfSeedOneOverallGain} and
\studyclaim{StudyTwoRfSeedTwoOverallGain}~MHz.  Their positive directions also
hold separately under interpolation and extrapolation
(with seed-level tables in the Supplementary Material). Correctness changes from
\studyclaim{StudyTwoSftCorrectness}\% for SFT to
\studyclaim{StudyTwoRfCorrectness}\% for RF. The primary endpoint accounts for
this change by assigning zero to every oracle-rejected draw. At the family
level, effects are positive for FIR, reverse FIR, polynomial, and recursive
filters, and negative for median filters, which appear only in extrapolation.
All leave-one-family-out mean effects remain positive, from
\studyclaim{StudyTwoSensitivityLeaveFamilyOutMin} to
\studyclaim{StudyTwoSensitivityLeaveFamilyOutMax}~MHz; detailed cells and
bootstrap sensitivities appear in the Supplementary Material.

The endpoint decomposition helps explain the gain. RF increases
pooled conditional frequency from \revisionclaim{SftConditional} to
\revisionclaim{RfConditional}~MHz despite its lower oracle pass rate.
Correctness-only reaches \revisionclaim{CorrectnessConditional}~MHz conditional
frequency. The improvement therefore comes from the physical quality of
passing draws. Fig.~\ref{fig:study2-controls} makes this distinction explicit:
correctness-only training reaches
\studyclaim{StudyTwoCorrectnessOnlyCorrectness}\% correctness, above both SFT
and RF, but attains only
\studyclaim{StudyTwoOverallCorrectnessOnlyPenalized}~MHz penalized \fmax.  Its
paired SFT-relative effect is
\studyclaim{StudyTwoOverallCorrectnessOnlySftGain}~MHz with interval
[\studyclaim{StudyTwoOverallCorrectnessOnlySftCiLow},
\studyclaim{StudyTwoOverallCorrectnessOnlySftCiHigh}]~MHz. A higher oracle pass
rate alone therefore does not explain the RF gain. The additional seeds
reported below test whether this contrast persists across training runs.

The original MLP reward does improve over SFT, reaching
\studyclaim{StudyTwoOverallMlpPenalized}~MHz with a paired gain of
\studyclaim{StudyTwoOverallMlpSftGain}~MHz.  The repaired RF nonetheless exceeds
it directly by \studyclaim{StudyTwoOverallRfMlpGain}~MHz with interval
[\studyclaim{StudyTwoOverallRfMlpCiLow},
\studyclaim{StudyTwoOverallRfMlpCiHigh}]~MHz.  The RF-minus-MLP intervals remain
above zero under interpolation and extrapolation as well.  Because the arms
share their SFT initialization, tasks, update budget, oracle, and evaluation,
this comparison favors the RF reward package for these checkpoints.

\begin{figure*}[t]
\centering
\includegraphics[width=.76\textwidth]{fig03_study2_controls}
\caption{Matched controls.  Panel (a) plots correctness against the
failure-penalized endpoint; the correctness-only arm is the key counterexample.
Panel (b) gives paired contrasts.  Repaired RF exceeds the original MLP reward
by \studyclaim{StudyTwoOverallRfMlpGain}~MHz
[\studyclaim{StudyTwoOverallRfMlpCiLow},
\studyclaim{StudyTwoOverallRfMlpCiHigh}]~MHz.}
\label{fig:study2-controls}
\end{figure*}
\label{sec:seed-extension}

% AUTO-GENERATED by paper/make_revision_evidence.py.
\begin{table}[t]
\centering
\caption{Completed post-primary setup-closure extension. Utility and paired differences are in MHz; intervals are exploratory \revisionclaim{SeedConfidence}\% design-bootstrap intervals conditional on checkpoints. Fail: distinct physical failures.}
\label{tab:seed-extension}
\small
\begin{tabular}{@{}lrrr@{}}
\toprule
Policy & Utility & Difference vs. SFT [CI] & Fail \\
\midrule
SFT & \revisionclaim{SeedSftMean} & Reference & \revisionclaim{SeedSftFailures} \\
RF seed three & \revisionclaim{SeedRfThreeMean} & \revisionclaim{SeedRfThreeGain} [\revisionclaim{SeedRfThreeLow}, \revisionclaim{SeedRfThreeHigh}] & \revisionclaim{SeedRfThreeFailures} \\
RF seed four & \revisionclaim{SeedRfFourMean} & \revisionclaim{SeedRfFourGain} [\revisionclaim{SeedRfFourLow}, \revisionclaim{SeedRfFourHigh}] & \revisionclaim{SeedRfFourFailures} \\
Correctness seed two & \revisionclaim{SeedCorrTwoMean} & \revisionclaim{SeedCorrTwoGain} [\revisionclaim{SeedCorrTwoLow}, \revisionclaim{SeedCorrTwoHigh}] & \revisionclaim{SeedCorrTwoFailures} \\
Correctness seed three & \revisionclaim{SeedCorrThreeMean} & \revisionclaim{SeedCorrThreeGain} [\revisionclaim{SeedCorrThreeLow}, \revisionclaim{SeedCorrThreeHigh}] & \revisionclaim{SeedCorrThreeFailures} \\
\bottomrule
\end{tabular}

\end{table}


The post-primary extension in Table~\ref{tab:seed-extension} compares the
additional seeds with SFT under the same setup-closure flow, retaining the
\revisionclaim{SeedFailureCount} physical failures as zero-valued outcomes.
RF reaches \revisionclaim{SeedRfThreeMean} and
\revisionclaim{SeedRfFourMean}~MHz against SFT's
\revisionclaim{SeedSftMean}~MHz.  Correctness-only reaches
\revisionclaim{SeedCorrTwoMean} and \revisionclaim{SeedCorrThreeMean}~MHz.
Both RF paired intervals are above zero and both correctness-only intervals
are below zero (Fig.~\ref{fig:seed-extension}), repeating the directions of
the original comparison. These runs share the original SFT initialization,
and their physical flow differs from the sealed primary flow. Their endpoints
and exploratory intervals are therefore reported separately from the primary
result.

\begin{figure}[t]
\centering
\includegraphics[width=\columnwidth]{fig07_seed_extension}
\caption{Separate post-primary seed extension. Points are paired per-design
utility differences; bars are exploratory
\revisionclaim{SeedConfidence}\% family-by-regime design-bootstrap intervals.
They condition on the evaluated checkpoints and sampled draws.}
\label{fig:seed-extension}
\end{figure}
\subsection{Physical costs and validation}
% AUTO-GENERATED by analyze_study2_secondary.py. DO NOT EDIT.
\begin{table*}[t]
\centering
\caption{Conditional physical-cost diagnostics on the
\studyclaim{StudyTwoPpaCommonSupport}-design three-arm common support. Values
are multiplicity-weighted within design and equally weighted across designs.
Power is Vivado's vectorless estimate, not board power.}
\label{tab:study2-ppa}
\small
\begin{tabular}{lrrrr}
\toprule
Metric & SFT & MLP & RF & RF--SFT difference [CI] \\
\midrule
LUTs & \studyclaim{StudyTwoPpaSftLut} & \studyclaim{StudyTwoPpaMlpLut} & \studyclaim{StudyTwoPpaRfLut} & \studyclaim{StudyTwoPpaRfSftLutDiff} [\studyclaim{StudyTwoPpaRfSftLutCiLow}, \studyclaim{StudyTwoPpaRfSftLutCiHigh}] \\
Flip-flops & \studyclaim{StudyTwoPpaSftFf} & \studyclaim{StudyTwoPpaMlpFf} & \studyclaim{StudyTwoPpaRfFf} & \studyclaim{StudyTwoPpaRfSftFfDiff} [\studyclaim{StudyTwoPpaRfSftFfCiLow}, \studyclaim{StudyTwoPpaRfSftFfCiHigh}] \\
DSP blocks & \studyclaim{StudyTwoPpaSftDsp} & \studyclaim{StudyTwoPpaMlpDsp} & \studyclaim{StudyTwoPpaRfDsp} & \studyclaim{StudyTwoPpaRfSftDspDiff} [\studyclaim{StudyTwoPpaRfSftDspCiLow}, \studyclaim{StudyTwoPpaRfSftDspCiHigh}] \\
BRAM blocks & \studyclaim{StudyTwoPpaSftBram} & \studyclaim{StudyTwoPpaMlpBram} & \studyclaim{StudyTwoPpaRfBram} & \studyclaim{StudyTwoPpaRfSftBramDiff} [\studyclaim{StudyTwoPpaRfSftBramCiLow}, \studyclaim{StudyTwoPpaRfSftBramCiHigh}] \\
Vectorless power [W] & \studyclaim{StudyTwoPpaSftPower} & \studyclaim{StudyTwoPpaMlpPower} & \studyclaim{StudyTwoPpaRfPower} & \studyclaim{StudyTwoPpaRfSftPowerDiff} [\studyclaim{StudyTwoPpaRfSftPowerCiLow}, \studyclaim{StudyTwoPpaRfSftPowerCiHigh}] \\
\bottomrule
\end{tabular}
\end{table*}


Higher frequency comes with a resource trade-off. On the
\studyclaim{StudyTwoPpaCommonSupport}-design common support, RF uses
\studyclaim{StudyTwoPpaRfLut} rather than
\studyclaim{StudyTwoPpaSftLut} LUTs and
\studyclaim{StudyTwoPpaRfFf} rather than
\studyclaim{StudyTwoPpaSftFf} flip-flops.  Conversely, DSP use falls from
\studyclaim{StudyTwoPpaSftDsp} to
\studyclaim{StudyTwoPpaRfDsp} blocks, while BRAM remains unchanged.
Vectorless power rises from \studyclaim{StudyTwoPpaSftPower} to
\studyclaim{StudyTwoPpaRfPower}~W.  The paired intervals in
Table~\ref{tab:study2-ppa} show that the LUT, register, DSP, and vectorless-power
directions are stable on common support.  These conditional quantities answer
what a compiled correct implementation costs; they do not replace the
failure-penalized policy endpoint.

% AUTO-GENERATED by analyze_study2_secondary.py. DO NOT EDIT.
\begin{table*}[t]
\centering
\caption{Illustrative source-level contrast for
\studyclaim{StudyTwoCaseDesign}. The entries are the dominant SFT candidate and
dominant candidate from RF seed one for that design. Selection occurred after
manuscript assessment; this example explains a possible CAD mechanism but is not a
population or causal estimate.}
\label{tab:study2-rtl-case}
\small
\begin{tabularx}{\textwidth}{lXrrrrr}
\toprule
Arm & Observed RTL organization & Sample mass [\%] & Fmax [MHz] & LUT & FF & DSP \\
\midrule
SFT & Delayed samples and product registers followed by a parallel final sum &
\studyclaim{StudyTwoCaseSftMass} & \studyclaim{StudyTwoCaseSftFmax} &
\studyclaim{StudyTwoCaseSftLut} & \studyclaim{StudyTwoCaseSftFf} &
\studyclaim{StudyTwoCaseSftDsp} \\
RF seed one & Registered transposed multiply--accumulate chain &
\studyclaim{StudyTwoCaseRfMass} & \studyclaim{StudyTwoCaseRfFmax} &
\studyclaim{StudyTwoCaseRfLut} & \studyclaim{StudyTwoCaseRfFf} &
\studyclaim{StudyTwoCaseRfDsp} \\
\bottomrule
\end{tabularx}
\end{table*}


The RTL pair in Table~\ref{tab:study2-rtl-case} illustrates this trade-off.  For \studyclaim{StudyTwoCaseDesign}, the dominant SFT
candidate stores delayed samples and products before a wide parallel sum.  The
dominant candidate from RF seed one instead emits a registered transposed
multiply--accumulate chain.  Vivado maps the latter to substantially more LUTs
and registers, no DSP blocks, and a much shorter timing path.  Because this
case was selected post-primary for interpretability, it illustrates
a physically meaningful RTL transformation but does not establish that every
family or every gain follows the same transformation.

Physical sensitivity checks support relative agreement with limits on absolute
frequency. The selected setup-closure pilot gives Spearman correlation
\revisionclaim{ClosureRho}, median symmetric error
\revisionclaim{ClosureMedianSmape}\%, and no reversal among
\revisionclaim{ClosureNonTiedPairs} non-tied pairs. Its worst candidate error is
\revisionclaim{ClosureWorstSmape}\%, on reverse FIR. Separately, the cross-version
coverage diagnostic confirms a setup path and clear constraint counts for
\revisionclaim{CoverageValid} of \revisionclaim{CoverageCandidates} candidates;
\revisionclaim{CoverageInvalid} lacks a path and
\revisionclaim{CoverageInconclusive} remain inconclusive. On common valid
candidates, median symmetric errors are \revisionclaim{CoverageSftMedianError}\%
for SFT, \revisionclaim{CoverageRfOneMedianError}\% for RF seed one, and
\revisionclaim{CoverageRfTwoMedianError}\% for RF seed two. The Supplementary
Material gives all cases and preserves the distinction between diagnostic
agreement and a revised primary endpoint.

The paired board measurements provide a further physical check
(Table~\ref{tab:study2-board} and Fig.~\ref{fig:board-evidence}). Across
\boardclaim{BoardRunCount} complete sweeps, the maximum
within-entry boundary spread is \boardclaim{BoardMaximumSpread}~MHz.  The echo
canary reaches \boardclaim{BoardCanaryFmax}~MHz and remains at least
\boardclaim{BoardMinimumCanaryMargin}~MHz above every DUT boundary. The
descriptive mean changes from \boardclaim{BoardSftMean} to
\boardclaim{BoardRfMean}~MHz, a \boardclaim{BoardMeanDifference}-MHz paired
difference driven by two large feed-forward gains. Polynomial and IIR tie,
while the median pair changes by \boardclaim{BoardMedDifference}~MHz. The
board results thus support the overall physical direction but show substantial
variation across families. As a descriptive subset measured under the amended
protocol in the Supplementary Material, they do not independently confirm the
population estimate.

\begin{figure}[t]
\centering
\includegraphics[width=.86\columnwidth]{fig04b_board_outcomes}
\caption{Descriptive PYNQ-Z2 clock boundaries under the amended protocol.
All pairs are retained: \boardclaim{BoardRfWins} RF wins,
\boardclaim{BoardTies} ties, and \boardclaim{BoardRfLosses} loss.
Annotations are signed RF-minus-SFT differences in MHz; the dashed line is
the echo-canary boundary. No population interval is inferred from this subset.}
\label{fig:board-evidence}
\end{figure}

% AUTO-GENERATED by analyze_study2_board.py. DO NOT EDIT.
\begin{table}[t]
\centering
\caption{Descriptive PYNQ-Z2 clock-sweep confirmation on
\boardclaim{BoardPairCount} paired designs. Values are medians over
\boardclaim{BoardRunCount} repeated sweeps. The subset required a disclosed
correctness-feasibility amendment and is not independent confirmatory evidence.}
\label{tab:study2-board}
\small
\begin{tabular}{llrrr}
\toprule
Family & Design & SFT & RF & RF--SFT [MHz] \\
\midrule
FIR & fir34\_v1\_8b & \boardclaim{BoardFirSft} & \boardclaim{BoardFirRf} & \boardclaim{BoardFirDifference} \\
FIRR & firr20\_v6 & \boardclaim{BoardFirrSft} & \boardclaim{BoardFirrRf} & \boardclaim{BoardFirrDifference} \\
POLY & poly16\_v13\_8b & \boardclaim{BoardPolySft} & \boardclaim{BoardPolyRf} & \boardclaim{BoardPolyDifference} \\
IIR & iir18\_v7 & \boardclaim{BoardIirSft} & \boardclaim{BoardIirRf} & \boardclaim{BoardIirDifference} \\
MED & med19 & \boardclaim{BoardMedSft} & \boardclaim{BoardMedRf} & \boardclaim{BoardMedDifference} \\
\bottomrule
\end{tabular}
\end{table}


\subsection{Structural changes and sampling efficiency}
\label{sec:mechanism}

The sampled candidates suggest that optimization often increases the
probability of implementations already accessible to SFT. Among
\studyclaim{StudyTwoSealedMechanismCommonCount} designs with passing sampled
strings under both policies, SFT's measured support contains a
candidate at least as fast as the RF conditional mean on
\studyclaim{StudyTwoSealedMechanismSupportCount}.  Mean inverse-Simpson
effective correct-string support contracts from
\studyclaim{StudyTwoSealedMechanismSftEffective} to
\studyclaim{StudyTwoSealedMechanismRfEffective}. Together, these observations
are consistent with a shift toward faster implementations on most designs
with common support. The samples are finite, however, and some sealed designs
fall outside this pattern, leaving circuit novelty unresolved.

On the sealed distribution, one RF draw lies between the empirical SFT
perfect-selector expectations for \revisionclaim{SelectorLowerN} draws
(\revisionclaim{SelectorLowerFmax}~MHz) and
\revisionclaim{SelectorUpperN} draws
(\revisionclaim{SelectorUpperFmax}~MHz).  At the upper integer crossing, the
SFT workload is \revisionclaim{SelectorUpperN} generations and
\revisionclaim{SelectorStreams} oracle simulation streams per design, with
\revisionclaim{SelectorRawImplementations} expected correct-draw implementation
attempts without caching or \revisionclaim{SelectorCachedImplementations}
distinct implementations with content caching.  These are empirical expected
workloads, not measured runtime or speedup.  They exclude policy-training cost
and assume perfect post-route selection on observed support; they do not imply
dominance over an unrestricted SFT search.

Source-feature aggregates give a complementary view across
\revisionclaim{StructureDesigns} designs with passing support in SFT, RF, and
correctness-only. The mean maximum multiplications per statement falls from
\revisionclaim{StructureStatementMultSft} for SFT to
\revisionclaim{StructureStatementMultRf} for RF, while correctness-only reaches
\revisionclaim{StructureStatementMultCorrectness}. RF also has higher structural
accumulation and pipeline-ratio scores. These descriptive lexical differences
are consistent with staged arithmetic and the measured FIR example. Because
some features also enter the reward, they are not independent causal evidence
or measurements of physical pipeline depth. The Supplementary Material reports
the complete comparison and its weighting.

\subsection{Backbone transfer and general RTL competence}
% AUTO-GENERATED by analyze_main_results.py. DO NOT EDIT.
\begin{table}[t]
\centering
\caption{Second-backbone replication on the earlier three-family split. Values
are penalized equal-sample post-route frequency; failures score zero.}
\label{tab:qwen-replication}
\small
\begin{tabular}{lrrrr}
\toprule
Regime & Designs & SFT & GRPO & Improved \\
\midrule
Interpolation & \claim{QwenInterpDesignCount} &
\claim{QwenInterpSftPenalized} & \claim{QwenInterpGrpoPenalized} &
\claim{QwenInterpImproved}/\claim{QwenInterpDesignCount} \\
Extrapolation & \claim{QwenExtrapDesignCount} &
\claim{QwenExtrapSftPenalized} & \claim{QwenExtrapGrpoPenalized} &
\claim{QwenExtrapImproved}/\claim{QwenExtrapDesignCount} \\
Overall & \claim{QwenDesignCount} & \claim{QwenOverallSftPenalized} &
\claim{QwenOverallGrpoPenalized} &
\claim{QwenOverallImproved}/\claim{QwenDesignCount} \\
\bottomrule
\end{tabular}
\end{table}

The full training recipe also transfers to Qwen-Coder on the earlier
\claim{QwenDesignCount}-design, \claim{QwenFamilyCount}-family universe.
Penalized \fmax{} rises from \claim{QwenOverallSftPenalized} to
\claim{QwenOverallGrpoPenalized}~MHz, a paired gain of
\claim{QwenOverallGain}~MHz with interval
[\claim{QwenOverallCiLow}, \claim{QwenOverallCiHigh}]~MHz; all
\claim{QwenDesignCount} designs improve.  Correctness also rises from
\claim{QwenOverallSftCorrectness}\% to
\claim{QwenOverallGrpoCorrectness}\%.  This supports transfer of the training recipe across backbones on the
earlier family universe. Its reward configuration and split differ from the
sealed RF experiment, so the two estimates remain separate.
% AUTO-GENERATED by paper/make_revision_evidence.py.
\begin{table}[t]
\centering
\caption{Matched post-primary general-RTL evaluation. Each checkpoint uses the same benchmark, generation settings, and draw budget.}
\label{tab:matched-verilog}
\small
\begin{tabular}{@{}lr@{}}
\toprule
Checkpoint & Pass-at-one (\%) \\
\midrule
SFT & \revisionclaim{VerilogSft} \\
RF seed one & \revisionclaim{VerilogRfOne} \\
RF seed two & \revisionclaim{VerilogRfTwo} \\
\bottomrule
\end{tabular}

\end{table}


The matched VerilogEval extension compares the actual sealed checkpoints on
\revisionclaim{VerilogProblems} official problems with
\revisionclaim{VerilogSamples} draws per checkpoint and problem.  Pass-at-one
is \revisionclaim{VerilogSft}\% for SFT,
\revisionclaim{VerilogRfOne}\% for RF seed one, and
\revisionclaim{VerilogRfTwo}\% for RF seed two
(Table~\ref{tab:matched-verilog}).  The RF-seed-mean paired change is
\revisionclaim{VerilogDelta} percentage points with interval
[\revisionclaim{VerilogCiLow}, \revisionclaim{VerilogCiHigh}], indicating a
small additional regression on general RTL tasks.
The interval resamples problems conditional on the evaluated checkpoints and
does not estimate policy-training seed variance.

Within the accelerator domain, the RF checkpoint trajectory retains the
physical gain from midpoint to endpoint. Interpolation changes from
\studyclaim{StudyTwoInterpMidFmax} to
\studyclaim{StudyTwoInterpEndFmax}~MHz and extrapolation from
\studyclaim{StudyTwoExtrapMidFmax} to
\studyclaim{StudyTwoExtrapEndFmax}~MHz. The Supplementary Material presents this
trajectory alongside the earlier reward failure, keeping the evidence from
the repaired run separate from the history that motivated it. Together with
the matched VerilogEval result, these checks show why physical gains within
the target domain and broader RTL competence should be evaluated separately.
The Supplementary Material also provides a runnable FIR example and a completed
RF versus SFT comparison on fresh draws with bounded routed setup-period
measurements. These post-primary results retain their own measurement contract
and are reported separately from the sealed primary estimate.
